\documentclass{article}
\usepackage[T1]{fontenc}
\usepackage{lmodern}
\usepackage{graphicx}
\usepackage{amsmath,amssymb}
\usepackage[a4paper,margin=2cm]{geometry}
\usepackage[absolute,overlay]{textpos}
\setlength{\TPHorizModule}{1mm}
\setlength{\TPVertModule}{1mm}
\usepackage[indonesian]{babel}
\usepackage{hyperref}
\setlength{\emergencystretch}{2em}
% unit: Halaman 4

\begin{document}

% === Halaman 4 (llm) ===

\section*{Kriptanalisis Cipher Abjad-Tunggal}

\begin{itemize}
    \item Kelemahan \textit{cipher} abjad-tunggal: tidak dapat menyembunyikan hubungan statistik antara plainteks dengan cipherteks.
    \begin{itemize}
        \item Huruf yang sama dienkripsi menjadi huruf cipherteks yang sama
        \item Huruf yang sering muncul di dalam plainteks, juga sering muncul di dalam huruf cipherteksnya yang berkoresponden.
        \item artinya struktur statistik bahasa masih muncul di dalam cipherteks
    \end{itemize}
    \item Oleh karena itu, cipherteks dapat didekripsi tanpa mengetahui kuncinya sekalipun dengan menggunakan \textbf{teknik analisis frekuensi}.
\end{itemize}

\end{document}
